\section{Introduction}

Conformal prediction~\cite{vovk2005algorithmic,shafer2008tutorial,angelopoulos2021gentle}
provides a general framework for attaching statistically valid uncertainty
quantification to black-box predictive models.  Given a model with predicted outputs~$A$ and true outcomes~$B$, conformal methods calibrate a distance
or nonconformity score~$f(A,B)$ on a held-out calibration set and choose a
threshold~$d$ so that, under exchangeability of samples, $
  \mathbb{P}(f(A^*,B^*)\le d) \ge \phi ,$ where $(A^*,B^*)$ corresponds to a new test pair.  The resulting
\emph{conformal set} $
  S(A^*) = \{B:\,f(A^*,B)\le d\}
$
contains the true outcome~$B^*$ with probability at least~$\phi$.
Conformal prediction thus converts any point predictor $A$ into a distribution-free set predictor $S(A)$ with finite-sample marginal coverage.


As an example, in  route planning and navigation systems, a predictive model proposes a route based on observed contextual features. Such features may include the user profile (e.g., commuter, tourist), real-time traffic conditions, weather, time of day, or the nature of the query (such as ``fastest'', ``scenic'', or ``kid-friendly'').  Despite the model's best prediction, a user may choose an alternate path for unmodeled reasons: personal preferences, temporary closures, or new information encountered en route.  Consequently, the model output $A_t$ (the suggested path) at a given epoch $t$ may differ from the actual path $B_t$ that the user traverses.  Over time, the system accumulates a history of such pairs
$\{(A_t,B_t)\}$, where $A_t$ is the model's suggested route and $B_t$ is the corresponding ground-truth route taken by the user.

The conformal prediction framework now works as follows:  For a new prediction $A^*$, we wish to construct a set of plausible user routes $S(A^*)$ that is guaranteed (under exchangeability) to contain the eventual ground-truth route $B^*$ with a prescribed probability level~$\phi$. 
Ideally we would also like $S(A^*)$ to be as compact or small as possible. 

%This conformal prediction fr provide interpretable, compressed uncertainty regions that reduce communication costs. 
%\emph{AV note to self: need to reword para}
The route planning task described above is an example of a conformal prediction task over graphs. 
%Such a framework over graphs generalizes to domains where realized behaviors diverge from model predictions. 
This framework also captures other settings where realized behaviors diverge from model predictions, for instance in \emph{trip planning}, where a model suggests an itinerary $A_t$ but the user follows $B_t$. Route planning also arises in routing autonomous vehicles, or in shipment logistics, while trip planning is akin to content recommendations. Other applications include settings where the model predicts %where the model recommends 
a sequence of clinical interventions (while physicians or patients may modify the plan), or predicts actions in business or manufacturing processes. These different applications can be modeled using a common graph-based conformal framework. 

%such as clinical treatment planning where the model suggests a treatment plan while the clinician chooses a different one; workflow recommendation, and inventory optimization. In each case, the conformal subgraph summarizes the space of likely outcomes, such as valid treatment pathways or utilized supply nodes, into a compact structure that guarantees coverage while minimizing complexity.


\subsection{Graph-based Conformal Compression} 

In structured prediction tasks such as route planning, itinerary recommendation, and logistics, the conformal set~$S(A^*)$ often becomes unwieldy, containing an enormous number of candidate outputs. Directly returning all elements of~$S(A^*)$, such as a raw, unstructured list of thousands of potential trajectories, is impractical and leads to cognitive overload for end users. We address this by introducing \emph{graph-based conformal compression}, a framework that summarizes the conformal set into a compact subgraph, preserving statistical coverage guarantees while minimizing structural complexity. 

Unlike an explicit enumeration of paths (which can grow exponentially), a \emph{conformal subgraph} provides a tractable and interpretable summary, even when the underlying predictive model is highly uncertain. Conversely, this framework provides a novel systems motivation for conformal prediction itself: by using past data to identify a high-probability, compressed subgraph, we safely prune the network and drastically reduce the search space for downstream routing, planning, and optimization algorithms.

Our main contribution is therefore in bridging the gap between conformal prediction and combinatorial optimization: We formally define the conformal subgraph problem and present a simple, computationally efficient algorithm that provably achieves (marginal) coverage in the distribution-free setting, while also being approximately efficient in terms of the size of the subgraph returned for that coverage bound.

In more detail, our main contributions are threefold:

\textbf{1. Algorithmic Formulation (Sections~\ref{sec:model},~\ref{sec:statement})} 
We develop an algorithmic, graph-based formulation of conformal compression that treats compression as the problem of selecting a compact subgraph  that contains the ground-truth output~$B^*$ with high probability.  We formalize the interaction between the conformal and compression guarantees in Section~\ref{sec:model}, and map the statistical goal of conformal coverage to a weighted hypergraph optimization problem in Section~\ref{sec:statement}. Our framework is flexible, accommodating (i) the classic distribution-free setting~\cite{vovk2005algorithmic} with variable graph contexts (i.e., $(A_t,B_t)$ for different $t$ corresponding to different underlying graphs and source-destination pairs), (ii) its specialization to a fixed graph context (i.e., the graph and source-destination pair are the same for different $t$), (iii) model probabilities derived directly from the predictive model \cite{angelopoulos2022learntestcalibratingpredictive,zhang2025conformalstructuredprediction} (where we assume access to a possibly mis-specified conditional distribution on $B$ given $A$), and (iv) risk-controlling prediction sets~\cite{Risk} to control general loss functions (e.g., controlling the expected fraction of missed route segments, or some other measure of deviation) rather than just 0-1 marginal coverage. We also show how our ideas can be applied to achieve conditional coverage~\cite{pmlr-v25-vovk12} in addition to marginal coverage.

\textbf{2. Efficiency Guarantees (Section~\ref{sec:algorithm}).} 
We connect conformal prediction with the classical optimization problem of \emph{densest-$k$-subgraph} selection in hypergraphs, where the goal is to select the subgraph with fewest vertices that preserves a specified fraction of the total weight of hyper-edges for a suitably defined weight function. For the resulting optimization problem, that we term the {\em conformal subgraph problem}, we develop provable guarantees on {\em efficiency}, that is, the size of the subgraph versus the conformal guarantee. We identify that conformal compression operates in a specific algorithmic regime where the subgraph retains a large fraction of hyper-edges, that is distinct from the classical, hard-to-approximate setting~\cite{bhaskara2010dks,feige2001dense,manurangsi2017hardness}. We show that in this high-coverage regime, the problem admits an efficient bicriteria constant-factor approximation (approximating both the vertex and hyper-edge set sizes) via a simple algorithm based on linear program (LP) rounding. We show in Appendix~\ref{sec:sampling} that this approach scales to exponentially large hyper-edge sets via sampling. 

\textbf{3. Statistical Validity via Monotonicity (Section~\ref{sec:algorithm}).} 
For a compression algorithm to support valid calibration in a distribution-free setting (or in the presence of a loss function~\cite{Risk}), it must produce a \emph{nested} sequence of subgraphs as the target coverage varies (monotonicity). Though the optimally compressed solution for a coverage bound is not monotone in the coverage bound (see Example~\ref{eg:monotone}), we prove that our approximation algorithm satisfies monotonicity. We show this by reinterpreting our algorithm via \emph{parametric minimum cuts}~\cite{gallo1989fast,ChekuriQT}, which also has the advantage of yielding an algorithm with quadratic running time. This result is one of our main contributions, and ensures that our method is a statistically valid conformal procedure for the classic distribution-free setting, both with variable and fixed graph contexts, while also enabling rigorous guarantees on both compression and validity. 

We view our paper as being primarily a theoretical paper, with the main contributions being conceptual or theoretical (in establishing provable validity and efficiency guarantees for conformal compression). Nevertheless, we present a set of experiments that provide a proof of concept that the proposed method is likely to work well in practice.  In Appendix~\ref{sec:experiments}, we complement our theoretical results with a simulation study on the canonical applications of trip planning and noisy navigation. In the former setting, vertices are activities, hyperedges are itineraries, and a conformal subgraph is a set of activities that captures a large fraction of itineraries; while in the latter setting, vertices are graph edges, hyperedges are navigation paths, and a conformal subgraph is a subgraph that captures a large fraction of routes. In both cases, we show that our algorithm either beats natural greedy baselines on the trade-off between compression and conformal calibration, or recovers a planted solution. We complement this with experiments on a real-world traffic dataset.


\subsection{Related Work}

\paragraph{Conformal Prediction.}
Conformal prediction has been applied across a wide range of areas
including regression~\cite{lei2018distribution,romano2019conformalized},
classification~\cite{vovk2005algorithmic,shafer2008tutorial},
structured models~\cite{zhang2025conformalstructuredprediction}, and more recently, outputs of language models~\cite{mohri2024languagemodelsconformalfactuality,cherian2024largelanguagemodelvalidity}. Our work differs in that it seeks compact representations of conformal prediction sets themselves, rather than improving the base predictor.   

There are recent extensions of conformal prediction to large or combinatorial decision problems, such as off-policy prediction~\cite{zhang2023conformal_offpolicy,taufiq2022conformal_copp} and reinforcement learning~\cite{plancp_2024,strawn2023conformal_psf}, though these works typically focus on constructing valid predictive intervals or uncertainty envelopes rather than on compressing such sets into small structural summaries.  %For example, several papers develop conformal off-policy prediction techniques that produce finite-sample predictive intervals for the return or outcomes of a target policy using logged data collected under a behavior policy~\cite{zhang2023conformal_offpolicy,taufiq2022conformal_copp}.  In reinforcement learning and control, conformal ideas have been applied to trajectory prediction and to design predictive safety filters that wrap learned controllers with uncertainty-aware safety checks providing high-probability coverage guarantees~\cite{plancp_2024,strawn2023conformal_psf}.  These works demonstrate the utility of conformal calibration in sequential decision settings, but their goal is to provide valid (often time-indexed) uncertainty sets or policy evaluation intervals, and they do not address the compression problem we study.

\paragraph{Efficiency in Conformal Prediction.} 
While conformal prediction guarantees validity by definition, minimizing the size (volume) of the prediction set, referred to as \textit{efficiency}, is challenging, with very  few recent results focusing on provable (as opposed to heuristic) efficiency of the solution (in our case, the size-optimality of the compressed graph). However, both provable validity and efficiency are a requirement, since validity is often trivial without efficiency (e.g., just output the entire graph).  %Guarantees for volume optimality typically rely on estimating the ``Highest Predictive Density'' (HPD) region. 

Works focusing on the dual goals of conformal validity and efficiency include \cite{lei2013distribution,sadinle2016least,izbicki2020flexible,izbicki2022cd,kiyani2024length} for prediction sets and regression. The works closest to our work are \cite{gao2025volume, gao2025highdimensional, gao2026ellipsoid} which provide conformal guarantees with approximate size optimality in an unsupervised setting for certain families of confidence sets like balls and ellipsoids. This is related to the fixed-context setting considered in Sections~\ref{sec:fixed} and~\ref{sec:fixed_optimal}.
%and~\cite{chernozhukov2021distributional,gao2025volume} for conditional guarantees. 
%As summarized in \cite{angelopoulos2024conformal}, a sufficient condition for asymptotic efficiency is the consistent estimation of the conditional density of the ground truth. However, these results require the latter space to be simple, e.g., low-dimensional reals or few classes. 
In contrast, in the combinatorial tasks we study, the output space is large and also precludes any posterior estimation. For instance, in routing, different calibration samples could have different source-destination pairs and different network edges. Our main contribution is developing techniques for conformal prediction in this more challenging setting that addresses both computational and statistical efficiency.

%This renders direct density estimation from limited samples statistically intractable. Consequently, we adopt a perspective similar to the above cited works, where validity is guaranteed for any estimator, while efficiency is achieved if the estimated quantiles converge to the truth. Similarly, our framework always guarantees validity via monotonicity (see Definition~\ref{def:monotone}), while we analyze efficiency (approximately optimal compression) in the ``perfectly calibrated'' regime where there is access to an approximately correct posterior. This formulation thus separates the statistical difficulty of learning from the combinatorial difficulty of extracting the optimal subgraph, and we explain it in more detail in Section~\ref{sec:model}.

Our work is also related to recent work on structured conformal prediction \cite{zhang2025conformalstructuredprediction}, which formulates conformal compression via model probabilities in the  ``learn--then--test'' framework~\cite{angelopoulos2022learntestcalibratingpredictive}. That work constructs interpretable prediction sets by searching within a structured label family specified as a DAG. In contrast, we start from a calibrated conformal set, which may be large and unstructured, and develop algorithmic methods for compressing it into a small subgraph. More crucially, their approach is not monotone and hence needs model probabilities, while we obtain distribution-free calibration via showing monotonicity of the approximate compressor. % The two approaches are  complementary: their machinery can be used  when the feasible family is structured as a DAG and interpretable, while our methods are applicable when the conformal set is large and arbitrarily structured. 



\paragraph{Densest Subgraph Problems.}
Our algorithms and analysis draw on a large body of work on densest subgraph discovery. The \emph{densest ratio subgraph} problem that maximizes the average degree admits efficient polynomial-time algorithms~\cite{goldberg1984finding,charikar2000greedy}, while its fixed-size variant, the \emph{densest-$k$-subgraph (DkS)} problem that restricts the subgraph to have at most $k$ vertices, is NP-hard and is conjectured to be inapproximable within polynomial factors~\cite{feige2001dense,bhaskara2010dks,manurangsi2017hardness}. The \emph{smallest-$p$-edge subgraph} (SpES) or \emph{minimum-$p$-union} problem~\cite{chlamtac2016dks} is a complementary minimization form, with similar inapproximability results.

The strong inapproximability results for densest $k$-subgraph are proved in parameter regimes where the target subgraph has number of edges sublinear in $m$, the total number of edges.  Those lower bounds therefore do not directly preclude efficient algorithms when the sought subgraph has $\Theta(m)$ edges, which is our regime of interest in conformal compression.  Further, the Lagrangian of our LP formulation in Section~\ref{sec:algorithm} is identical to the Lagrangian for densest ratio subhypergraphs, hence reducing the problem to \emph{parametric min-cut} on bipartite graphs~\cite{goldberg1984finding,gallo1989fast,ChekuriQT}. This connection is crucial to showing monotonicity and valid calibration. %, and as far as we are aware, this is not present in existing literature. Our results therefore fill a gap between the polynomial-time solvable unconstrained density problems and the small-subgraph hardness regime emphasized in prior work.

%\paragraph{Overview.}
%Section~\ref{sec:model} introduces our modeling framework and formalizes  conformal compression. Section~\ref{sec:statement}  develops a hypergraph densest subgraph abstraction that we call the {\em conformal subgraph problem}.  Section~\ref{sec:algorithm} presents an algorithm for this problem with a constant factor bicriteria approximation guarantee, and that satisfies monotonicity and hence conformal validity. Appendix~\ref{app:np-hard} shows {\sc NP-Hardness} in the regime of interest, and Appendix~\ref{sec:sampling} shows how to handle the setting with exponentially many hyper-edges. We present experimental results in Section~\ref{sec:experiments} for synthetic and real-world trip planning and navigation scenarios, showing that our algorithm handily beats simple greedy baselines. 
